\section{Implementation}
\label{sec:implementation}

The Linux prototype turns the policy contract into a real
\code{cgroup/namei_ext} attachment path. Policies are eBPF programs under
\file{bpf/policies/*.bpf.c}. The implementation consists of the VFS call sites,
the BPF ABI, user-space policy attachment and map setup, and kernel-side
validation for the supported action set.

\subsection{Path-Walk Integration}

The prototype integrates policy at three VFS points. \code{walk_component()}
handles intermediate and ordinary trailing components,
\code{open_last_lookups()} handles the final component for open/create, and
\code{iterate_dir()} mediates entries returned by a lower directory iterator.
Each call site first tests the global enable key and parent-scope filter. A
matching lookup constructs the common context and invokes the same
\code{cgroup/namei_ext} decision function before the VFS commits to a lower
result.

The component and final-open wrappers preserve the surrounding namei state
rather than starting a second path walk. \code{PASS} enters the ordinary
component or final-open implementation. \code{REDIRECT} replaces only the
current \code{qstr}, invokes the lower directory's \code{d_hash()} operation,
and restores the original component after lookup. \code{HIDE} returns
\code{ENOENT}. \code{SELECT_TARGET} installs a validated registered path into
the existing \code{nameidata}; intermediate walking or final VFS completion
then continues from that object.

Cache-hot lookup reaches these wrappers in RCU-walk. The BPF-visible flags mask
the internal \code{LOOKUP_RCU} bit so cache state does not change the policy
input. \code{PASS} remains in RCU-walk. Same-parent redirect converts through
\code{try_to_unlazy()} before changing the component, while errors preserve
the VFS conversion and \code{RESOLVE_CACHED} behavior. The wrapper does not
reinvoke the program merely to apply a saved action after successful
conversion. If conversion returns \code{ECHILD}, the surrounding VFS may
restart the complete pathname walk and invoke policy again. Programs therefore
cannot assume exactly-once execution across VFS restart and must make
observable side effects idempotent.

\subsection{BPF ABI}

The kernel-facing ABI exposes one decision function. The context includes an
event type, operation flags, the current component name, parent identity,
identifiers used by the policy, and bounded output fields for the selected
path-view action. The current action set is \code{PASS}, \code{REDIRECT},
\code{HIDE}, and \code{SELECT_TARGET}. \code{REDIRECT} carries a bounded
replacement component in \code{redirect_name}. \code{SELECT_TARGET} carries an
opaque \code{target_id} that indexes a kernel-registered lower \code{struct
path}. The policy cannot return arbitrary unbounded path strings or perform
recursive path walks.

The verifier exposes names and parent identity as read-only context and permits
writes only to the bounded replacement-name and target-ID outputs. After
program execution, the kernel validates action type, name length and contents,
target identity, object type, and selection scope before continuing. Redirect
decisions reject empty names, \code{.}, \code{..}, embedded nulls, slashes, and
create operations.

\subsection{Registered-Target Lifetime}

Workload controllers open each target with \code{O_PATH} and register its file
descriptor under an opaque target ID. The kernel copies the descriptor's
\code{struct path}, holds mount and dentry references, and publishes the record
in an RCU hash table keyed by the policy scope and ID. Replacing or clearing a
record uses two related paths: replacement atomically publishes the new RCU
hash entry, while clearing removes the old entry. Both wait for an RCU grace
period before dropping the retired path references. A policy never receives
the path or a kernel pointer.

Ref-walk target selection looks up the record and acquires independent path
references. RCU-walk instead borrows the registry's references while the outer
namei RCU critical section protects the record. When namei installs a borrowed
target, it samples the target dentry sequence and retains the mount sequence.
Existing \code{try_to_unlazy()} and \code{complete_walk()} validation either
obtain stable references or reject and restart a changed walk. Thus concurrent
target replacement cannot free a path still borrowed by an RCU reader.

Before installation, the kernel rejects missing IDs, symbolic links,
unsupported object types, create-through-select, scoped walks, and prohibited
cross-mount selection. Intermediate targets must be directories. A final target
may be an existing regular file or directory, after which ordinary stat,
access, open, exec, \code{O_PATH}, or \code{O_DIRECTORY} handling proceeds on
the selected object. The prototype does not create targets, copy data, populate
caches, or implement lifecycle orchestration.

\subsection{Replacement-Component Validation And Permission Preservation}

Workload setup configures the policy state that chooses replacement components.
The prototype implements the redirect subset by copying a bounded
\code{redirect_name} from the BPF output, validating its length, action, and
same-parent scope, and then resuming the normal VFS path. The replacement
component is not an
arbitrary path string and does not trigger a recursive path walk. After
validation, the normal VFS path continues, including permission checks and
lower-filesystem file operations. Same-parent validation keeps \namei from
becoming arbitrary path rewriting or a filesystem service. Both component
replacement and registered-target selection reject \code{LOOKUP_CREATE}: they
may choose an existing object, but they cannot change lower-filesystem creation
or copy-up rules.

The implemented hide action uses the same decision point but does not allocate
or synthesize filesystem objects. Lookup observes absence, and directory
enumeration suppresses the lower entry for the active policy state.

\subsection{Directory Iteration}

\code{iterate_dir()} wraps the lower filesystem's existing
\code{dir_context.actor}. For every returned lower entry, the wrapper invokes
the common decision function. \code{PASS} forwards the original name and inode
information, \code{REDIRECT} forwards a validated replacement name, and
\code{HIDE} consumes the lower entry without exposing it. The wrapper copies
the iterator position, count, and directory-entry flags in both directions so
partial \code{getdents64()} calls retain the lower iterator's progress.

\code{SELECT_TARGET} is intentionally rejected as a per-entry readdir action.
Selecting an existing directory during lookup and then enumerating that opened
directory is supported; inventing a new alias for an arbitrary target is not.
The latter would require synthetic-entry storage and stable position semantics
that the current existing-object boundary does not own.

\subsection{Unused And Unmanaged Fast Paths}

The cgroup BPF attachment uses a static key for the globally disabled case.
With no active program, component lookup, final-open lookup, and directory
iteration enter their ordinary implementations after that test. An attached
program may additionally register up to eight exact parent paths. An RCU hash
index and active-global counter let all three call sites reject an unrelated
parent before entering a noinline wrapper, constructing a BPF context, or
discovering the current cgroup.

Scope updates run under the scope mutex, while a raw spinlock sequence counter
publishes the global counters, RCU hash nodes, and active state without a
sleeping lock in RCU-walk. Removed scope records and held path references are
released only after an RCU grace period. The prefilter is conservative: a
positive result still executes the effective-owner and per-policy scope checks,
while counter overflow saturates toward policy execution rather than creating a
false negative.

\subsection{Failure Handling}

Verifier-invalid programs cannot attach. At runtime, an unknown action,
malformed redirect, missing registration, incompatible target type, or
unsupported operation returns a specific error to the pathname operation. The
implementation does not translate those errors into \code{PASS}; doing so
would expose the ordinary lower name when the requested policy did not execute.
This visible error behavior and the unchanged lower object after invalid
decisions are implementation-correctness checks; RQ3 primarily compares which
filesystem methods and runtime state each mechanism owns.

Phase 1 validation runs in KVM with the modified kernel and the real
\code{cgroup/namei_ext} attach path. Host-only checks, object inspection, and
bpftool-only smoke tests are diagnostics rather than validation for the paper's
mechanism claims.

These components give the evaluation a single validation path in which Make
targets boot the modified kernel, attach policies through
\code{cgroup/namei_ext}, and preserve raw validation and experiment artifacts.
